\documentclass[a4paper,10pt]{report}\input{../0.Préambule/Préambule_cours_prof}\begin{document}

\newpage
\setcounter{chapter}{9}
\chapter{Proportionnalité 1}

\section{Tableau de proportionnalité}

\begin{dft}
Dans un tableau de nombres à deux lignes, on reconnaît une situation de proportionnalité lorsque les nombres
de la deuxième ligne s'obtiennent en multipliant ceux de la première par un même nombre.
\medskip

\noindent Ce nombre est appelé coefficient de proportionnalité.
\end{dft}

\begin{ex}
\begin{center}
\renewcommand{\arraystretch}{1.8}
\begin{tabularx}{12cm}{|>{\columncolor{lightgray!50}}p{3cm}|*{5}{>{\centering \arraybackslash}X|}}
\hline
Nombre de sms & $1$ & $2$ & $5$ & $8$ & $10$ \\\hline
Prix payé en \euro & $0,15$ & $0,30$ & $0,75$ & $1,20$ & $1,50$ \\\hline
\end{tabularx}
\end{center}

\medskip
\noindent On a : \quad $0,15 \div 1 = 0,15$ \quad ;\quad  $0,30 \div 2 = 0,15$ \quad ;\quad  $0,75 \div 5 = 0,15$ \quad ;\quad  $1,20 \div 8 = 0,15$ \quad et \quad $1,50 \div 10 = 0,15$
\medskip

\noindent Tous les quotients sont égaux à $0,15$ donc le prix payé est proportionnel au nombre de sms envoyés.
\medskip

\noindent Le coefficient de proportionnalité est $0,15$. Cela signifie qu'on paie $0,15$\euro pour $1$ sms.
\end{ex}

\begin{ex}

\begin{center}
\renewcommand{\arraystretch}{1.8}
\begin{tabularx}{9cm}{|>{\columncolor{lightgray!50}}p{3cm}|*{3}{>{\centering \arraybackslash}X|}}
\hline
Nombre de BD & $3$ & $5$ & $11$ \\\hline
Prix payé & $7,50$ & $9$ & $18$ \\\hline
\end{tabularx}
\end{center}
\medskip

\noindent On a : \quad $7,50 \div 3 = 2,5 \quad \quad 9 \div 5 = 1,8$
\medskip

\noindent $2,5 \neq 1,8$ donc le prix n'est pas proportionnel au nombre de BD achetées.

\noindent Ce tableau n'est pas un tableau de proportionnalité.
\end{ex}

\section{Quatrième proportionnelle}

\begin{regle}
Dans un tableau de proportionnalité, si on connaît $3$ nombres non nuls, (dont deux se correspondent), alors on
peut calculer le quatrième nombre manquant.
\medskip

\noindent Ce nombre manquant est appelé quatrième proportionnelle.
\end{regle}

\begin{meth}
Lorsqu'on achète des cerises, leur prix (en \euro) est proportionnel à leur masse (en kg).
\medskip

\noindent Voici différentes méthodes pour calculer le nombre manquant $x$ :

\begin{enumerate}
\item \textbf{On utilise le coefficient de proportionnalité :}
\begin{center}
\renewcommand{\arraystretch}{1.8}
\begin{tabularx}{7cm}{|>{\columncolor{lightgray!50}}p{3cm}|*{2}{>{\centering \arraybackslash}X|}}
\hline
Masse (en kg) & $4$ & $5$ \\\hline
Prix (en \euro) & $11,20$ & $x$ \\\hline
\end{tabularx}
\end{center}
Le coefficient de proportionnalité est $11,20 \div 4 = 2,8$

\noindent $5 \times  2,8 = 14$ donc $x = 14$

\newpage
\item \textbf{Méthode des produits en croix :}
\begin{center}
\renewcommand{\arraystretch}{1.8}
\begin{tabularx}{7cm}{|>{\columncolor{lightgray!50}}p{3cm}|*{2}{>{\centering \arraybackslash}X|}}
\hline
Masse (en kg) & $4$ & $5$ \\\hline
Prix (en \euro) & $11,20$ & $x$ \\\hline
\end{tabularx}
\end{center}
On a : 

\vspace{-5mm}
\begin{eqnarray*}
\dfrac{11,20}{4} &=& \dfrac{x}{5} \\
x &=& \dfrac{11,20 \times 5}{4}\\
x &=& \dfrac{56}{4}\\
x &=& 14
\end{eqnarray*}
On multiplie les nombres des cases en diagonale et on divise par le troisième nombre.
\item \textbf{Par multiplication (ou division) :}

\medskip
\noindent On peut multiplier ou diviser une colonne :

\vspace{5mm}
\begin{center}
\renewcommand{\arraystretch}{1.8}
\begin{tabularx}{7cm}{|>{\columncolor{lightgray!50}}p{3cm}|*{2}{>{\centering \arraybackslash}X|}}
\hline
Masse (en kg) & $4$ & $5$ \\\hline
Prix (en \euro) & $11,20$ & $x$ \\\hline
\end{tabularx}
\end{center}
\vspace{5mm}

$11,20 \times 1,25 = 14$ donc $x = 14$.
\item \textbf{Par addition :}

\medskip
\begin{center}
\renewcommand{\arraystretch}{1.8}
\begin{tabularx}{9cm}{|>{\columncolor{lightgray!50}}p{3cm}|*{3}{>{\centering \arraybackslash}X|}}
\hline
Masse (en kg) & $4$ & $1$ & $5$ \\\hline
Prix (en \euro) & $11,20$ & $2,80$ & $x$ \\\hline
\end{tabularx}
\end{center}

\medskip
\noindent On a : \quad $4 + 1 = 5$ donc $11,20 + 2, 80 = 14$. 

\noindent Donc $x = 14$.
\end{enumerate}
\end{meth}
\end{document}